\documentclass[12pt,a4paper]{article}

% Packages
\usepackage[utf8]{inputenc}
\usepackage[T1]{fontenc}
\usepackage{amsmath,amssymb,amsfonts}
\usepackage{mathtools}
\usepackage{graphicx}
\usepackage{booktabs}
\usepackage{hyperref}
\usepackage[margin=2.5cm]{geometry}
\usepackage{natbib}
\usepackage{xcolor}
\usepackage{multirow}
\usepackage{siunitx}
\usepackage{array}

% Custom commands
\newcommand{\hdecrate}{$h$\,decay rate}
\newcommand{\SEdecrate}{SampEn\,decay rate}
\newcommand{\AEdecrate}{ApEn\,decay rate}
\DeclareMathOperator{\MI}{MI}

\title{Input-Signal Robustness of Multiscale Entropy Biomarkers:\\
The Permutation Entropy Rate as a Uniquely Invariant\\
Complexity Measure for Cardiac Time Series}
\author{Martin G.\ Frasch$^{1}$ \and Nicolas B.\ Garnier$^{2}$}

\newcommand{\affiliations}{%
$^{1}$Department of Obstetrics and Gynecology, University of Washington, Seattle, WA, USA\\
$^{2}$Laboratoire de Physique, CNRS, ENS de Lyon, Universit\'{e} de Lyon, Lyon, France
}
\date{\today}

\begin{document}
\maketitle
\begingroup
\centering
\small
\affiliations
\par
\endgroup
\vspace{1em}

\begin{abstract}
Entropy-based complexity measures derived from cardiac interval time series are sensitive to seemingly arbitrary preprocessing choices: whether the input is heart rate (HR) or R-R interval (RRI), and whether the signal is normalized.
We systematically evaluate the robustness of three multiscale entropy decay rates---permutation entropy rate ($h$), Sample Entropy (SampEn), and Approximate Entropy (ApEn)---across four input conditions (\{HR, RRI\} $\times$ \{raw, normalized\}) in 251~segments from 49~children.
The $h$~decay rate shows near-perfect cross-condition agreement (Pearson $r = 0.983$--$0.998$) and maintains a significant dose-response relationship ($p < 0.006$, Cohen's $d > 1.5$) in all four conditions.
In contrast, SampEn and ApEn decay rates are highly sensitive to input choice ($r$ as low as $-0.06$), with statistical significance depending on the specific condition tested.
We provide a theoretical explanation grounded in the rank-based nature of k-nearest-neighbor estimation, which is invariant under the monotonic HR$\leftrightarrow$RRI transform, whereas the absolute tolerance threshold used in template matching is not.
These results establish the permutation entropy decay rate as the uniquely robust complexity biomarker for cardiac time series and caution against interpreting template-matching entropy results without explicit robustness verification.
\end{abstract}


\section{Introduction}\label{sec:intro}

Multiscale entropy analysis has become a standard tool for characterizing the complexity of physiological time series~\citep{costa2002,costa2005}.
In cardiac applications, entropy is typically computed from either heart-rate (HR, in beats per minute) or R-R interval (RRI, in milliseconds) time series---two representations of the same underlying process related by the reciprocal transform $\text{HR} = 60{,}000 / \text{RRI}$.
Additional choices include whether to normalize the signal by its standard deviation, which affects the absolute scale of inter-point distances.

While these preprocessing choices are often treated as inconsequential, they can profoundly affect entropy estimates.
Template-matching algorithms (Sample Entropy, Approximate Entropy) use an absolute tolerance threshold $r_{\text{abs}} = r \times \text{SD}(x)$ that depends on the distributional geometry of the input~\citep{richman2000}.
The nonlinear HR$\leftrightarrow$RRI transform alters this geometry, potentially changing which templates are counted as matches.

In contrast, k-nearest-neighbor (k-NN) entropy estimators~\citep{kozachenko1987,kraskov2004} operate on rank-ordered distances, a property that confers invariance under monotonic transforms.
This theoretical distinction has practical consequences that, to our knowledge, have not been systematically evaluated in the context of multiscale cardiac complexity analysis.

We address this gap by computing three families of multiscale entropy measures under four input conditions and evaluating their robustness at three levels: (i)~segment-level cross-condition correlations, (ii)~preservation of dose-response group differences, and (iii)~stability of effect sizes.
The dataset comprises ambulatory ECG recordings from a cohort study of prenatal steroid exposure~\citep{frasch2024}, providing a clinical context in which biomarker robustness directly affects interpretability.

A companion paper~\citep{garnier2026} describes the numerical validation of the Python entropy library used for the k-NN computations.


\section{Methods}\label{sec:methods}

\subsection{Entropy measures}\label{sec:entropy}

We compute three families of entropy from cardiac interval time series, each capturing different aspects of signal complexity.

\subsubsection{Permutation entropy rate ($h$)}

The permutation entropy rate quantifies the conditional uncertainty of the next value given the past, estimated via k-nearest-neighbor (k-NN) distances using the Kozachenko--Leonenko (KL) estimator~\citep{kozachenko1987}:
\begin{equation}\label{eq:kl}
  \hat{H}(X) = d \cdot \frac{1}{N_{\text{v}}} \sum_{\substack{i=1 \\ \varepsilon_i > 0}}^{N} \log(2\varepsilon_i)
  + \psi(N_{\text{v}}) - \psi(k),
\end{equation}
where $\varepsilon_i$ is the $L_\infty$ (Chebyshev) distance to the $k$-th nearest neighbor, $\psi(\cdot)$ is the digamma function, and $N_{\text{v}} \leq N$ is the number of points with $\varepsilon_i > 0$.

The entropy rate at time scale $\tau$ is computed via the decomposition~\citep{kraskov2004}:
\begin{equation}\label{eq:erate}
  \hat{h}(\tau) = \hat{H}(X_t) - \hat{I}(X_t;\, X_{\text{past}}^{(m)}),
\end{equation}
where $X_{\text{past}}^{(m)} = (X_{t-\tau}, X_{t-2\tau}, \ldots, X_{t-m\tau})$ and $\hat{I}$ is estimated via the KSG algorithm~\citep{kraskov2004}.
Both $\hat{H}$ and $\hat{I}$ are computed on the same unified time-delay embedding to ensure algebraic consistency (see~\citealt{garnier2026} for implementation details).

\subsubsection{Sample Entropy (SampEn)}

SampEn quantifies signal regularity through template matching~\citep{richman2000}:
\begin{equation}\label{eq:sampen}
  \text{SampEn}(m, r) = -\ln\!\left(\frac{B^{m+1}}{A^m}\right),
\end{equation}
where $A^m$ and $B^{m+1}$ count the number of template pairs matching within tolerance $r_{\text{abs}} = r \times \text{SD}(x)$ at embedding dimensions $m$ and $m+1$, respectively, using the $L_\infty$ norm and excluding self-matches.

\subsubsection{Approximate Entropy (ApEn)}

ApEn follows the same template-matching framework but includes self-matches in the counting~\citep{pincus1991}, making it less sensitive to short data but introducing a bias toward lower values.

\subsection{Multiscale analysis}\label{sec:multiscale}

\subsubsection{Scale sweep}

Entropy is computed at time scales ranging from 0.5~to 5.0~s in steps of $d\tau = 0.04$~s (i.e., \texttt{dstride}~$= 10$ samples at 250~Hz), yielding approximately 113~scale points per segment.

For SampEn and ApEn, multiscale analysis uses non-overlapping block averaging (coarse-graining;~\citealt{costa2002}):
\begin{equation}\label{eq:coarsegrain}
  y_j^{(\tau)} = \frac{1}{\tau} \sum_{i=(j-1)\tau+1}^{j\tau} x_i.
\end{equation}
For the entropy rate $h$, multiscale behavior is captured implicitly through the stride parameter $\tau$ in the time-delay embedding~(Eq.~\ref{eq:erate}), operating on the original signal with progressively wider lags rather than on a downsampled signal.

\subsubsection{Summary statistics}

From each entropy-versus-scale curve, four summary statistics are extracted:
\begin{itemize}
  \item \textbf{AUC}: area under the curve (trapezoidal integration)
  \item \textbf{max}: maximum entropy value across scales
  \item \textbf{max\_t}: scale at which the maximum occurs
  \item \textbf{dec\_rate}: rate of decrease from the maximum to the final scale
\end{itemize}
This yields 12~features per segment: \{$h$, SampEn, ApEn\} $\times$ \{AUC, max, max\_t, dec\_rate\}.

\subsection{Computation parameters}\label{sec:params}

\begin{table}[htbp]
\centering
\caption{Entropy computation parameters.}\label{tab:params}
\begin{tabular}{@{}lll@{}}
\toprule
Parameter & Value & Notes \\
\midrule
Embedding dimension $m$ & 2 & Entropy rate and SampEn/ApEn \\
k-NN neighbors $k$ & 5 & KL estimator \\
Tolerance $r$ (SampEn/ApEn) & 0.2 $\times$ SD & Fixed across scales \\
Distance metric & $L_\infty$ (Chebyshev) & Both k-NN and template matching \\
Scale range & 0.5--5.0~s & $d\tau = 0.04$~s (\texttt{dstride}${}=10$) \\
$N_{\text{eff}}$ & 4096 & Points per realization \\
$N_{\text{real}}$ & 3 & Realizations averaged \\
\bottomrule
\end{tabular}
\end{table}

\subsection{Input conditions}\label{sec:conditions}

All entropy measures were computed under four input-signal variants (Table~\ref{tab:conditions}), each with and without temporal smoothing (8~total configurations).

\begin{table}[htbp]
\centering
\caption{Four input conditions for robustness analysis.}\label{tab:conditions}
\begin{tabular}{@{}llll@{}}
\toprule
Label & Signal & Normalization & Transform from RRI \\
\midrule
HR\textsubscript{raw} & Heart rate (bpm) & None & $60{,}000 / \text{RRI}$ \\
HR\textsubscript{norm} & Heart rate & $\div$ SD & Monotonic \\
RRI\textsubscript{raw} & R-R intervals (ms) & None & Identity \\
RRI\textsubscript{norm} & R-R intervals & $\div$ SD & Monotonic \\
\bottomrule
\end{tabular}
\end{table}

\subsection{Dataset}\label{sec:dataset}

We analyze 251~five-minute ECG segments (250~Hz) from 49~children aged 8~years, drawn from a prospective cohort study of prenatal maternal steroid exposure~\citep{frasch2024}.
Subjects are classified into three groups based on cumulative maternal steroid dose: low~dose ($<5$~g, $N = 12$--15), high~dose ($\geq 5$~g, $N = 13$--15), and controls ($N = 24$--30), with exact counts varying by segment due to quality exclusions.
R-R intervals are extracted via ensemble R-peak detection with signal quality index (SQI) filtering.

\subsection{Statistical analysis}\label{sec:stats}

\subsubsection{Cross-condition correlations}
Segment-level Pearson correlations between decay rates computed under different input conditions.

\subsubsection{Dose-response preservation}
Kruskal-Wallis tests comparing three dose groups on the mental arithmetic segment (M3\_bis\_M4, peak cognitive load), repeated under each input condition.

\subsubsection{Pairwise effect sizes}
Mann-Whitney $U$ tests with Cohen's $d$ for the low-dose vs.\ high-dose comparison under each condition.


\section{Results}\label{sec:results}

\subsection{Cross-condition correlations}\label{sec:corr}

Table~\ref{tab:correlations} shows segment-level Pearson correlations ($N = 251$) between decay rates computed under the four input conditions.

\begin{table}[htbp]
\centering
\caption{Cross-condition Pearson correlations for entropy decay rates.
The $h$~decay rate shows near-perfect agreement ($r > 0.98$) across all conditions.
SampEn and ApEn decay rates are highly sensitive to input choice.}\label{tab:correlations}
\small
\begin{tabular}{l cccccc}
\toprule
& \multicolumn{2}{c}{$h$ decay rate} & \multicolumn{2}{c}{SampEn decay rate} & \multicolumn{2}{c}{ApEn decay rate} \\
\cmidrule(lr){2-3} \cmidrule(lr){4-5} \cmidrule(lr){6-7}
Pair & $r$ & & $r$ & & $r$ & \\
\midrule
HR\textsubscript{raw} vs HR\textsubscript{norm}   & \textbf{0.997} & & 0.713 & & 0.091 \\
HR\textsubscript{raw} vs RRI\textsubscript{raw}   & \textbf{0.983} & & 0.502 & & 0.099 \\
HR\textsubscript{raw} vs RRI\textsubscript{norm}  & \textbf{0.985} & & 0.698 & & 0.089 \\
HR\textsubscript{norm} vs RRI\textsubscript{raw}  & \textbf{0.985} & & 0.316 & & $-0.045$ \\
HR\textsubscript{norm} vs RRI\textsubscript{norm} & \textbf{0.986} & & 0.953 & & 0.980 \\
RRI\textsubscript{raw} vs RRI\textsubscript{norm} & \textbf{0.998} & & 0.310 & & $-0.060$ \\
\midrule
Range & \textbf{0.983--0.998} & & 0.310--0.953 & & $-0.060$--0.980 \\
\bottomrule
\end{tabular}
\end{table}

The $h$~decay rate is virtually invariant to input choice, with all pairwise correlations exceeding 0.98.
In contrast, SampEn decay rate correlations range from 0.31 to 0.95, and ApEn decay rate correlations range from $-0.06$ to 0.98.
For both template-matching measures, only the HR\textsubscript{norm} vs RRI\textsubscript{norm} pair achieves reasonable agreement ($r > 0.95$), while comparisons involving raw signals show moderate to negligible correlations.

\subsection{Dose-response preservation}\label{sec:kw}

Table~\ref{tab:kw} presents Kruskal-Wallis tests comparing three dose groups on the mental arithmetic segment under each input condition.

\begin{table}[htbp]
\centering
\caption{Dose-response analysis (Kruskal-Wallis) on the M3\_bis\_M4 segment across input conditions.
Only the $h$~decay rate reaches significance ($p < 0.01$) in all four conditions.}\label{tab:kw}
\small
\begin{tabular}{l cc cc cc}
\toprule
& \multicolumn{2}{c}{$h$ decay rate} & \multicolumn{2}{c}{SampEn decay rate} & \multicolumn{2}{c}{ApEn decay rate} \\
\cmidrule(lr){2-3} \cmidrule(lr){4-5} \cmidrule(lr){6-7}
Condition & $H$ & $p$ & $H$ & $p$ & $H$ & $p$ \\
\midrule
HR\textsubscript{raw}  & 10.16 & \textbf{0.006} & 12.49 & \textbf{0.002} & 1.90 & 0.387 \\
HR\textsubscript{norm} & 10.59 & \textbf{0.005} & 4.09  & 0.129           & 4.41 & 0.110 \\
RRI\textsubscript{raw} & 10.18 & \textbf{0.006} & 4.51  & 0.105           & 0.37 & 0.832 \\
RRI\textsubscript{norm}& 10.69 & \textbf{0.005} & 4.49  & 0.106           & 4.51 & 0.105 \\
\bottomrule
\end{tabular}
\end{table}

The $h$~decay rate yields consistent $H$ statistics (10.16--10.69) and significant $p$~values ($p < 0.006$) across all four conditions.
The SampEn decay rate is significant only under HR\textsubscript{raw} ($p = 0.002$) and fails to reach significance in all other conditions.
The ApEn decay rate is not significant in any condition.

\subsection{Pairwise comparisons}\label{sec:pairwise}

Table~\ref{tab:pairwise} shows Mann-Whitney $U$ tests for the low-dose vs.\ high-dose comparison.

\begin{table}[htbp]
\centering
\caption{Low-dose vs.\ high-dose pairwise comparisons (Mann-Whitney $U$)
on the M3\_bis\_M4 segment.
The $h$~decay rate shows consistently large effects ($d > 1.5$, $p < 0.005$) in all conditions.}\label{tab:pairwise}
\small
\begin{tabular}{l ccc ccc}
\toprule
& \multicolumn{3}{c}{$h$ decay rate} & \multicolumn{3}{c}{SampEn decay rate} \\
\cmidrule(lr){2-4} \cmidrule(lr){5-7}
Condition & $U$ & $p$ & $d$ & $U$ & $p$ & $d$ \\
\midrule
HR\textsubscript{raw}  & 25 & \textbf{0.004} & 1.53 & 14 & \textbf{$<$0.001} & 1.97 \\
HR\textsubscript{norm} & 23 & \textbf{0.003} & 1.57 & 42 & 0.054               & 0.93 \\
RRI\textsubscript{raw} & 25 & \textbf{0.004} & 1.57 & 21 & 0.062               & 0.97 \\
RRI\textsubscript{norm}& 24 & \textbf{0.004} & 1.66 & 41 & 0.047               & 0.98 \\
\bottomrule
\end{tabular}
\end{table}

The $h$~decay rate shows large, stable effects (Cohen's $d = 1.53$--$1.66$) across all conditions, with $p < 0.005$ uniformly.
The SampEn decay rate achieves the largest single effect size ($d = 1.97$) under HR\textsubscript{raw}, but this effect does not generalize: in the remaining three conditions, the $p$~values hover near $0.05$ and the effect sizes drop to $d \approx 0.95$.

\subsection{Summary of robustness properties}\label{sec:summary}

\begin{table}[htbp]
\centering
\caption{Summary of robustness properties across entropy decay rates.}\label{tab:summary}
\begin{tabular}{@{}l ccc@{}}
\toprule
Property & $h$ dec.\ rate & SampEn dec.\ rate & ApEn dec.\ rate \\
\midrule
Cross-condition $r$ (min) & \textbf{0.983} & 0.310 & $-0.060$ \\
Significant in all 4 (KW) & \textbf{Yes (4/4)} & No (1/4) & No (0/4) \\
Signif.\ low vs.\ high, all 4 (MW) & \textbf{Yes (4/4)} & No (1--2/4) & No (0--1/4) \\
Effect size range ($d$) & 1.53--1.66 & 0.93--1.97 & 0.05--0.93 \\
HR\textsubscript{norm} vs RRI\textsubscript{norm} $r$ & \textbf{0.986} & 0.953 & 0.980 \\
\bottomrule
\end{tabular}
\end{table}


\section{Discussion}\label{sec:discussion}

\subsection{Theoretical basis: k-NN rank invariance}\label{sec:theory}

The robustness of the $h$~decay rate follows from the rank-based nature of the k-NN estimator.
The KL estimator (Eq.~\ref{eq:kl}) depends on neighbor distances $\varepsilon_i$ through their logarithms.
Under a monotonic transform $g(x)$:
\begin{enumerate}
  \item The \emph{rank ordering} of inter-point distances is preserved---the same points remain nearest neighbors.
  \item The distances scale by the local derivative $|g'|$, adding a term $\overline{\ln|g'|}$ to the entropy estimate.
  \item Entropy \emph{differences} across scales---which determine the decay rate---are \emph{invariant}, because the constant offset cancels in the subtraction.
\end{enumerate}

The HR$\leftrightarrow$RRI conversion ($g(x) = 60{,}000/x$) is a strictly monotonic (decreasing) transform.
Normalization by standard deviation ($x \mapsto x/\sigma$) is also monotonic.
Therefore, while absolute entropy values differ across conditions (by a constant $\overline{\ln|g'|}$), the \emph{shape} of the entropy-versus-scale curve---and hence its decay rate---is preserved.
This explains the cross-condition correlations $r > 0.98$ observed for the $h$~decay rate.

\subsection{Template-matching sensitivity}\label{sec:template}

SampEn and ApEn use an absolute tolerance threshold $r_{\text{abs}} = 0.2 \times \text{SD}(x)$.
Under the nonlinear HR$\leftrightarrow$RRI transform:
\begin{enumerate}
  \item The standard deviation changes nonlinearly: $\text{SD}(60{,}000/x) \neq 60{,}000 \cdot \text{SD}(x) / \overline{x}^2$ in general.
  \item The tolerance $r_{\text{abs}}$ captures a \emph{different proportion} of the signal's distributional geometry.
  \item The set of matching templates changes, sometimes dramatically.
\end{enumerate}

This effect is asymmetric across the HR/RRI range: at low HR values (long RRI), a given RRI tolerance corresponds to a small HR tolerance, while at high HR values the same RRI tolerance corresponds to a large HR tolerance.
The low cross-condition correlations for SampEn ($r$ as low as 0.31) and ApEn ($r$ as low as $-0.06$) directly reflect this sensitivity.

\subsection{Normalization as partial remedy}\label{sec:normalization}

Normalization by standard deviation partially mitigates the template-matching sensitivity.
The HR\textsubscript{norm} vs RRI\textsubscript{norm} correlations are high for all three measures ($r = 0.953$--$0.986$; Table~\ref{tab:correlations}, last row).
However, this agreement only holds between normalized variants; comparisons involving raw signals remain unstable.
Moreover, normalization does not restore the dose-response relationship: even under RRI\textsubscript{norm}, only the $h$~decay rate retains significance (Table~\ref{tab:kw}).

This observation has important practical implications: while normalization improves reproducibility of template-matching measures, it does not make them robust biomarkers.
The $h$~decay rate, by contrast, yields the same statistical conclusion regardless of preprocessing choice.

\subsection{The SampEn ``HR\textsubscript{raw} anomaly''}\label{sec:anomaly}

The SampEn decay rate achieves its strongest result ($H = 12.49$, $p = 0.002$, $d = 1.97$) specifically under HR\textsubscript{raw}---exceeding even the $h$~decay rate in that single condition.
This isolated finding does not replicate under any other input condition.
We interpret it as an artifact of the particular distributional geometry of raw heart rate: the reciprocal transform compresses long RRI values (low HR) and expands short RRI values (high HR), creating a specific template-matching landscape that happens to amplify group differences.
Had only HR\textsubscript{raw} been tested, SampEn decay rate would have appeared to be the stronger biomarker---a cautionary example of condition-dependent findings.

\subsection{Relation to implementation validation}\label{sec:validation}

A companion paper~\citep{garnier2026} validates the Python implementation of the k-NN entropy estimators against the reference C/ANN library, identifying and correcting three implementation-level discrepancies (zero-distance handling, distance-metric mismatch, embedding alignment).
The present work complements that validation by demonstrating that k-NN entropy measures are not only numerically correct but also \emph{methodologically robust}---their biological conclusions are invariant to preprocessing choices that fundamentally alter template-matching results.

The two papers share a common insight: the rank-based nature of k-NN estimation is the key property.
In~\citet{garnier2026}, it ensures that the corrected implementation produces identical neighbor rankings as the C reference.
Here, it ensures that the same neighbor rankings---and hence the same entropy decay rates---emerge regardless of whether the input is HR or RRI, raw or normalized.

\subsection{Limitations}\label{sec:limitations}

This analysis evaluates robustness across input-signal transformations but does not address other sources of variability such as segment length, R-peak detection method, or noise contamination.
The dose-response analysis is limited to a single cognitive-stress segment (mental arithmetic) in a modest cohort ($N = 49$), and replication in larger datasets is needed.
The theoretical rank-invariance argument is exact only for strictly monotonic transforms; nonlinear transforms that change the rank order (e.g., absolute value) would break the invariance.


\section{Conclusion}\label{sec:conclusion}

Among three families of multiscale entropy decay rates, only the permutation entropy rate ($h$~decay rate) is robust to the choice of input signal representation.
Its cross-condition correlations exceed $r = 0.98$, and it maintains a large, significant dose-response effect ($d > 1.5$, $p < 0.005$) under all four tested conditions.
The theoretical basis is the rank invariance of k-NN distance estimation under monotonic transforms---a property absent from the template-matching framework used by SampEn and ApEn.

We recommend that studies reporting entropy-based cardiac complexity measures:
\begin{enumerate}
  \item Use k-NN-based entropy rates as the primary complexity biomarker.
  \item Verify robustness across input conditions before interpreting template-matching results.
  \item Report the specific input condition (HR vs.\ RRI, raw vs.\ normalized) to enable reproducibility.
\end{enumerate}

All analysis code and data are available at \url{https://github.com/martinfrasch/entropy} and \url{https://github.com/martinfrasch/ECG_steroids}.


\section*{Acknowledgments}

The robustness analysis and manuscript preparation were performed with assistance from Claude~Code (Anthropic).


\bibliographystyle{abbrvnat}
\bibliography{references_robustness}

\end{document}
